% Shared preamble for CS 377 worksheets.
%
% Usage: a worksheet begins with % Shared preamble for CS 377 worksheets.
%
% Usage: a worksheet begins with % Shared preamble for CS 377 worksheets.
%
% Usage: a worksheet begins with \input{00-template.tex}, sets its metadata with
% \wscourse / \wstitle, then opens \begin{document} and calls
% \wsheader (group version) or \wsheaderindividual (individual version).
%
% Build with build.sh (pdflatex).

\documentclass[11pt]{article}

\usepackage[T1]{fontenc}
\usepackage[utf8]{inputenc}
\usepackage{lmodern}
\usepackage[margin=0.9in, top=0.75in, bottom=0.75in]{geometry}
\usepackage{microtype}
\usepackage{array}
\usepackage{tabularx}
\usepackage{calc}
\usepackage{enumitem}
\usepackage{forloop}

\setlength{\parindent}{0pt}
\setlength{\parskip}{0.5\baselineskip}
\pagestyle{empty}
\raggedbottom

% Vertically center the contents of every X column, so that short entries sit in
% the middle of a tall ruled row rather than clinging to its top edge.
\renewcommand{\tabularxcolumn}[1]{m{#1}}

% ---------------------------------------------------------------------------
% Metadata
% ---------------------------------------------------------------------------
\newcommand{\wscourseval}{CS 377}
\newcommand{\wstitleval}{Worksheet}

\newcommand{\wscourse}[1]{\renewcommand{\wscourseval}{#1}}
\newcommand{\wstitle}[1]{\renewcommand{\wstitleval}{#1}}

% Section meeting times offered in the current term, one per column of the
% "circle one" row. Fall 2026 has two sections. To add or drop one, edit both
% \wssections and \wssectioncount.
\newcommand{\wssections}{\wssectiontime{12:30pm} & \wssectiontime{2:00pm}}
\newcommand{\wssectioncount}{2}

% ---------------------------------------------------------------------------
% Header pieces
% ---------------------------------------------------------------------------

% Course number, small and right-aligned above the title.
\newcommand{\wscourseline}{%
  {\raggedleft\small\wscourseval\par}%
  \vspace{0.4\baselineskip}%
}

\newcommand{\wsbigtitle}{%
  {\fontsize{26}{30}\selectfont\bfseries\wstitleval\par}%
  \vspace{0.4\baselineskip}%
}

% Four ruled boxes (2x2) for group members to sign in.
\newcommand{\wsmembertable}{%
  \begin{tabularx}{\textwidth}{|X|X|}
    \hline
    \rule{0pt}{2.4\baselineskip} & \\
    \hline
    \rule{0pt}{2.4\baselineskip} & \\
    \hline
  \end{tabularx}%
}

% Single ruled box for a name.
\newcommand{\wsnametable}{%
  \begin{tabularx}{\textwidth}{|>{\raggedleft\arraybackslash\bfseries}m{1.4in}|X|}
    \hline
    Your Name & \rule{0pt}{2.4\baselineskip} \\
    \hline
  \end{tabularx}%
}

% A cell of fixed height whose contents are vertically centered.
%
% An m-column is not enough on its own: it centers the cell's *box*, which for a
% multi-line paragraph leaves the text riding high against the top rule. Giving
% every cell in the row an explicit equal-height parbox with [c] inner alignment
% centers the text itself, whether it is one line or two.
\newcommand{\wsvcell}[2]{\parbox[c][#1][c]{\linewidth}{#2}}

% "Section (circle one)" row with one column per meeting time. The row height is
% measured once at the surrounding font size, so that the \small label and the
% normal-size times agree on how tall the row is.
\newlength{\wssectionrowht}
\newcommand{\wssectiontime}[1]{\wsvcell{\wssectionrowht}{\centering #1}}
\newcommand{\wssectiontable}{%
  \setlength{\wssectionrowht}{2.6\baselineskip}%
  \begin{tabularx}{\textwidth}{|>{\bfseries\small}m{1.1in}|*{\wssectioncount}{X|}}
    \hline
    \wsvcell{\wssectionrowht}{\raggedleft Section\newline(circle one)} & \wssections \\
    \hline
  \end{tabularx}%
}

% Group worksheet header: course, title, group members, section.
\newcommand{\wsheader}{%
  \wscourseline
  \wsbigtitle
  Group members:\par
  \vspace{0.2\baselineskip}
  \wsmembertable
  \vspace{0.6\baselineskip}
  \wssectiontable
  \vspace{0.4\baselineskip}
}

% Individual worksheet header: course/date, title, name, section.
\newcommand{\wsheaderindividual}{%
  \wscourseline
  \wsbigtitle
  \wsnametable
  \vspace{0.6\baselineskip}
  \wssectiontable
  \vspace{0.4\baselineskip}
}

% ---------------------------------------------------------------------------
% Questions
% ---------------------------------------------------------------------------
\newcounter{wsquestion}

% \wsquestion{prompt}{writing space} -- auto-numbered prompt followed by blank
% space for students to write in.
\newcommand{\wsquestion}[2]{%
  \stepcounter{wsquestion}%
  \par\noindent(\thewsquestion)~#1\par
  \vspace{#2}%
}

% Prompt with no trailing space, for when a table or list follows.
\newcommand{\wsprompt}[1]{%
  \stepcounter{wsquestion}%
  \par\noindent(\thewsquestion)~#1\par
}

% Worksheet tables: a header row of italic headings over a stack of tall blank
% rows for students to write in.
%
% \wsgridtable{left heading}{right heading}{number of rows}{row height}
% \wsgridtablethree{left}{middle}{right}{number of rows}{row height}
%
% The blank rows are accumulated into \wsrowsbody *before* the tabularx is
% opened. Running the loop inside the table instead leaves a stray extra row,
% because the loop's own machinery counts as the start of one.
\newcounter{wsrow}
\newcounter{wscol}
\makeatletter
\newcommand{\wsrowsbody}{}
% \wsbuildrows{number of rows}{row height}{number of columns}
\newcommand{\wsbuildrows}[3]{%
  \gdef\wsrowsbody{}%
  \forloop{wsrow}{0}{\value{wsrow} < #1}{%
    \g@addto@macro\wsrowsbody{\rule[-0.5\baselineskip]{0pt}{#2}}%
    \forloop{wscol}{1}{\value{wscol} < #3}{%
      \g@addto@macro\wsrowsbody{&}%
    }%
    \g@addto@macro\wsrowsbody{\\ \hline}%
  }%
}
\makeatother

% Each heading gets an equal-height centered cell, so that the headings line up
% with each other and sit centered in the heading row rather than against its
% top rule.
\newlength{\wsheadht}
\newcommand{\wshead}[1]{\wsvcell{\wsheadht}{\itshape #1}}

\newcommand{\wsgridtable}[4]{%
  \wsbuildrows{#3}{#4}{2}%
  \setlength{\wsheadht}{1.6\baselineskip}%
  \begin{tabularx}{\textwidth}{|X|X|}
    \hline
    \wshead{#1} & \wshead{#2} \\
    \hline
    \wsrowsbody
  \end{tabularx}%
}

\newcommand{\wsgridtablethree}[5]{%
  \wsbuildrows{#4}{#5}{3}%
  \setlength{\wsheadht}{1.6\baselineskip}%
  \begin{tabularx}{\textwidth}{|X|X|X|}
    \hline
    \wshead{#1} & \wshead{#2} & \wshead{#3} \\
    \hline
    \wsrowsbody
  \end{tabularx}%
}
, sets its metadata with
% \wscourse / \wstitle, then opens \begin{document} and calls
% \wsheader (group version) or \wsheaderindividual (individual version).
%
% Build with build.sh (pdflatex).

\documentclass[11pt]{article}

\usepackage[T1]{fontenc}
\usepackage[utf8]{inputenc}
\usepackage{lmodern}
\usepackage[margin=0.9in, top=0.75in, bottom=0.75in]{geometry}
\usepackage{microtype}
\usepackage{array}
\usepackage{tabularx}
\usepackage{calc}
\usepackage{enumitem}
\usepackage{forloop}

\setlength{\parindent}{0pt}
\setlength{\parskip}{0.5\baselineskip}
\pagestyle{empty}
\raggedbottom

% Vertically center the contents of every X column, so that short entries sit in
% the middle of a tall ruled row rather than clinging to its top edge.
\renewcommand{\tabularxcolumn}[1]{m{#1}}

% ---------------------------------------------------------------------------
% Metadata
% ---------------------------------------------------------------------------
\newcommand{\wscourseval}{CS 377}
\newcommand{\wstitleval}{Worksheet}

\newcommand{\wscourse}[1]{\renewcommand{\wscourseval}{#1}}
\newcommand{\wstitle}[1]{\renewcommand{\wstitleval}{#1}}

% Section meeting times offered in the current term, one per column of the
% "circle one" row. Fall 2026 has two sections. To add or drop one, edit both
% \wssections and \wssectioncount.
\newcommand{\wssections}{\wssectiontime{12:30pm} & \wssectiontime{2:00pm}}
\newcommand{\wssectioncount}{2}

% ---------------------------------------------------------------------------
% Header pieces
% ---------------------------------------------------------------------------

% Course number, small and right-aligned above the title.
\newcommand{\wscourseline}{%
  {\raggedleft\small\wscourseval\par}%
  \vspace{0.4\baselineskip}%
}

\newcommand{\wsbigtitle}{%
  {\fontsize{26}{30}\selectfont\bfseries\wstitleval\par}%
  \vspace{0.4\baselineskip}%
}

% Four ruled boxes (2x2) for group members to sign in.
\newcommand{\wsmembertable}{%
  \begin{tabularx}{\textwidth}{|X|X|}
    \hline
    \rule{0pt}{2.4\baselineskip} & \\
    \hline
    \rule{0pt}{2.4\baselineskip} & \\
    \hline
  \end{tabularx}%
}

% Single ruled box for a name.
\newcommand{\wsnametable}{%
  \begin{tabularx}{\textwidth}{|>{\raggedleft\arraybackslash\bfseries}m{1.4in}|X|}
    \hline
    Your Name & \rule{0pt}{2.4\baselineskip} \\
    \hline
  \end{tabularx}%
}

% A cell of fixed height whose contents are vertically centered.
%
% An m-column is not enough on its own: it centers the cell's *box*, which for a
% multi-line paragraph leaves the text riding high against the top rule. Giving
% every cell in the row an explicit equal-height parbox with [c] inner alignment
% centers the text itself, whether it is one line or two.
\newcommand{\wsvcell}[2]{\parbox[c][#1][c]{\linewidth}{#2}}

% "Section (circle one)" row with one column per meeting time. The row height is
% measured once at the surrounding font size, so that the \small label and the
% normal-size times agree on how tall the row is.
\newlength{\wssectionrowht}
\newcommand{\wssectiontime}[1]{\wsvcell{\wssectionrowht}{\centering #1}}
\newcommand{\wssectiontable}{%
  \setlength{\wssectionrowht}{2.6\baselineskip}%
  \begin{tabularx}{\textwidth}{|>{\bfseries\small}m{1.1in}|*{\wssectioncount}{X|}}
    \hline
    \wsvcell{\wssectionrowht}{\raggedleft Section\newline(circle one)} & \wssections \\
    \hline
  \end{tabularx}%
}

% Group worksheet header: course, title, group members, section.
\newcommand{\wsheader}{%
  \wscourseline
  \wsbigtitle
  Group members:\par
  \vspace{0.2\baselineskip}
  \wsmembertable
  \vspace{0.6\baselineskip}
  \wssectiontable
  \vspace{0.4\baselineskip}
}

% Individual worksheet header: course/date, title, name, section.
\newcommand{\wsheaderindividual}{%
  \wscourseline
  \wsbigtitle
  \wsnametable
  \vspace{0.6\baselineskip}
  \wssectiontable
  \vspace{0.4\baselineskip}
}

% ---------------------------------------------------------------------------
% Questions
% ---------------------------------------------------------------------------
\newcounter{wsquestion}

% \wsquestion{prompt}{writing space} -- auto-numbered prompt followed by blank
% space for students to write in.
\newcommand{\wsquestion}[2]{%
  \stepcounter{wsquestion}%
  \par\noindent(\thewsquestion)~#1\par
  \vspace{#2}%
}

% Prompt with no trailing space, for when a table or list follows.
\newcommand{\wsprompt}[1]{%
  \stepcounter{wsquestion}%
  \par\noindent(\thewsquestion)~#1\par
}

% Worksheet tables: a header row of italic headings over a stack of tall blank
% rows for students to write in.
%
% \wsgridtable{left heading}{right heading}{number of rows}{row height}
% \wsgridtablethree{left}{middle}{right}{number of rows}{row height}
%
% The blank rows are accumulated into \wsrowsbody *before* the tabularx is
% opened. Running the loop inside the table instead leaves a stray extra row,
% because the loop's own machinery counts as the start of one.
\newcounter{wsrow}
\newcounter{wscol}
\makeatletter
\newcommand{\wsrowsbody}{}
% \wsbuildrows{number of rows}{row height}{number of columns}
\newcommand{\wsbuildrows}[3]{%
  \gdef\wsrowsbody{}%
  \forloop{wsrow}{0}{\value{wsrow} < #1}{%
    \g@addto@macro\wsrowsbody{\rule[-0.5\baselineskip]{0pt}{#2}}%
    \forloop{wscol}{1}{\value{wscol} < #3}{%
      \g@addto@macro\wsrowsbody{&}%
    }%
    \g@addto@macro\wsrowsbody{\\ \hline}%
  }%
}
\makeatother

% Each heading gets an equal-height centered cell, so that the headings line up
% with each other and sit centered in the heading row rather than against its
% top rule.
\newlength{\wsheadht}
\newcommand{\wshead}[1]{\wsvcell{\wsheadht}{\itshape #1}}

\newcommand{\wsgridtable}[4]{%
  \wsbuildrows{#3}{#4}{2}%
  \setlength{\wsheadht}{1.6\baselineskip}%
  \begin{tabularx}{\textwidth}{|X|X|}
    \hline
    \wshead{#1} & \wshead{#2} \\
    \hline
    \wsrowsbody
  \end{tabularx}%
}

\newcommand{\wsgridtablethree}[5]{%
  \wsbuildrows{#4}{#5}{3}%
  \setlength{\wsheadht}{1.6\baselineskip}%
  \begin{tabularx}{\textwidth}{|X|X|X|}
    \hline
    \wshead{#1} & \wshead{#2} & \wshead{#3} \\
    \hline
    \wsrowsbody
  \end{tabularx}%
}
, sets its metadata with
% \wscourse / \wstitle, then opens \begin{document} and calls
% \wsheader (group version) or \wsheaderindividual (individual version).
%
% Build with build.sh (pdflatex).

\documentclass[11pt]{article}

\usepackage[T1]{fontenc}
\usepackage[utf8]{inputenc}
\usepackage{lmodern}
\usepackage[margin=0.9in, top=0.75in, bottom=0.75in]{geometry}
\usepackage{microtype}
\usepackage{array}
\usepackage{tabularx}
\usepackage{calc}
\usepackage{enumitem}
\usepackage{forloop}

\setlength{\parindent}{0pt}
\setlength{\parskip}{0.5\baselineskip}
\pagestyle{empty}
\raggedbottom

% Vertically center the contents of every X column, so that short entries sit in
% the middle of a tall ruled row rather than clinging to its top edge.
\renewcommand{\tabularxcolumn}[1]{m{#1}}

% ---------------------------------------------------------------------------
% Metadata
% ---------------------------------------------------------------------------
\newcommand{\wscourseval}{CS 377}
\newcommand{\wstitleval}{Worksheet}

\newcommand{\wscourse}[1]{\renewcommand{\wscourseval}{#1}}
\newcommand{\wstitle}[1]{\renewcommand{\wstitleval}{#1}}

% Section meeting times offered in the current term, one per column of the
% "circle one" row. Fall 2026 has two sections. To add or drop one, edit both
% \wssections and \wssectioncount.
\newcommand{\wssections}{\wssectiontime{12:30pm} & \wssectiontime{2:00pm}}
\newcommand{\wssectioncount}{2}

% ---------------------------------------------------------------------------
% Header pieces
% ---------------------------------------------------------------------------

% Course number, small and right-aligned above the title.
\newcommand{\wscourseline}{%
  {\raggedleft\small\wscourseval\par}%
  \vspace{0.4\baselineskip}%
}

\newcommand{\wsbigtitle}{%
  {\fontsize{26}{30}\selectfont\bfseries\wstitleval\par}%
  \vspace{0.4\baselineskip}%
}

% Four ruled boxes (2x2) for group members to sign in.
\newcommand{\wsmembertable}{%
  \begin{tabularx}{\textwidth}{|X|X|}
    \hline
    \rule{0pt}{2.4\baselineskip} & \\
    \hline
    \rule{0pt}{2.4\baselineskip} & \\
    \hline
  \end{tabularx}%
}

% Single ruled box for a name.
\newcommand{\wsnametable}{%
  \begin{tabularx}{\textwidth}{|>{\raggedleft\arraybackslash\bfseries}m{1.4in}|X|}
    \hline
    Your Name & \rule{0pt}{2.4\baselineskip} \\
    \hline
  \end{tabularx}%
}

% A cell of fixed height whose contents are vertically centered.
%
% An m-column is not enough on its own: it centers the cell's *box*, which for a
% multi-line paragraph leaves the text riding high against the top rule. Giving
% every cell in the row an explicit equal-height parbox with [c] inner alignment
% centers the text itself, whether it is one line or two.
\newcommand{\wsvcell}[2]{\parbox[c][#1][c]{\linewidth}{#2}}

% "Section (circle one)" row with one column per meeting time. The row height is
% measured once at the surrounding font size, so that the \small label and the
% normal-size times agree on how tall the row is.
\newlength{\wssectionrowht}
\newcommand{\wssectiontime}[1]{\wsvcell{\wssectionrowht}{\centering #1}}
\newcommand{\wssectiontable}{%
  \setlength{\wssectionrowht}{2.6\baselineskip}%
  \begin{tabularx}{\textwidth}{|>{\bfseries\small}m{1.1in}|*{\wssectioncount}{X|}}
    \hline
    \wsvcell{\wssectionrowht}{\raggedleft Section\newline(circle one)} & \wssections \\
    \hline
  \end{tabularx}%
}

% Group worksheet header: course, title, group members, section.
\newcommand{\wsheader}{%
  \wscourseline
  \wsbigtitle
  Group members:\par
  \vspace{0.2\baselineskip}
  \wsmembertable
  \vspace{0.6\baselineskip}
  \wssectiontable
  \vspace{0.4\baselineskip}
}

% Individual worksheet header: course/date, title, name, section.
\newcommand{\wsheaderindividual}{%
  \wscourseline
  \wsbigtitle
  \wsnametable
  \vspace{0.6\baselineskip}
  \wssectiontable
  \vspace{0.4\baselineskip}
}

% ---------------------------------------------------------------------------
% Questions
% ---------------------------------------------------------------------------
\newcounter{wsquestion}

% \wsquestion{prompt}{writing space} -- auto-numbered prompt followed by blank
% space for students to write in.
\newcommand{\wsquestion}[2]{%
  \stepcounter{wsquestion}%
  \par\noindent(\thewsquestion)~#1\par
  \vspace{#2}%
}

% Prompt with no trailing space, for when a table or list follows.
\newcommand{\wsprompt}[1]{%
  \stepcounter{wsquestion}%
  \par\noindent(\thewsquestion)~#1\par
}

% Worksheet tables: a header row of italic headings over a stack of tall blank
% rows for students to write in.
%
% \wsgridtable{left heading}{right heading}{number of rows}{row height}
% \wsgridtablethree{left}{middle}{right}{number of rows}{row height}
%
% The blank rows are accumulated into \wsrowsbody *before* the tabularx is
% opened. Running the loop inside the table instead leaves a stray extra row,
% because the loop's own machinery counts as the start of one.
\newcounter{wsrow}
\newcounter{wscol}
\makeatletter
\newcommand{\wsrowsbody}{}
% \wsbuildrows{number of rows}{row height}{number of columns}
\newcommand{\wsbuildrows}[3]{%
  \gdef\wsrowsbody{}%
  \forloop{wsrow}{0}{\value{wsrow} < #1}{%
    \g@addto@macro\wsrowsbody{\rule[-0.5\baselineskip]{0pt}{#2}}%
    \forloop{wscol}{1}{\value{wscol} < #3}{%
      \g@addto@macro\wsrowsbody{&}%
    }%
    \g@addto@macro\wsrowsbody{\\ \hline}%
  }%
}
\makeatother

% Each heading gets an equal-height centered cell, so that the headings line up
% with each other and sit centered in the heading row rather than against its
% top rule.
\newlength{\wsheadht}
\newcommand{\wshead}[1]{\wsvcell{\wsheadht}{\itshape #1}}

\newcommand{\wsgridtable}[4]{%
  \wsbuildrows{#3}{#4}{2}%
  \setlength{\wsheadht}{1.6\baselineskip}%
  \begin{tabularx}{\textwidth}{|X|X|}
    \hline
    \wshead{#1} & \wshead{#2} \\
    \hline
    \wsrowsbody
  \end{tabularx}%
}

\newcommand{\wsgridtablethree}[5]{%
  \wsbuildrows{#4}{#5}{3}%
  \setlength{\wsheadht}{1.6\baselineskip}%
  \begin{tabularx}{\textwidth}{|X|X|X|}
    \hline
    \wshead{#1} & \wshead{#2} & \wshead{#3} \\
    \hline
    \wsrowsbody
  \end{tabularx}%
}


\wscourse{CS 377}
\wstitle{Virtue Ethics: Group Reflection}

\begin{document}

\wsheader

\textit{Work through these questions together as a group. Turn in one sheet per group at the end of class.}

\wsprompt{As a group, discuss your role models (i.e. ``exemplars''). Write down three of them below, along with the \textbf{virtues} your group thinks they emulate in their character traits. For this exercise, they must be real people.}

\vfill

\wsprompt{Now do the same for ``anti-role models'' (i.e. negative exemplars). Write down two or three, along with the \textbf{vices} your group sees in their character traits. Again, they must be real people.}

\vfill

\textit{(Questions 3 and 4 are on the back)}

\newpage

\wsprompt{According to virtue ethics, a person's happiness and flourishing (i.e. eudaemonia) is connected to morals and values: to flourish, a person must live in accordance with virtues. Write down character-focused virtues you think a typical UIC student might wish to develop this semester, then discuss as a group and write down a concrete way that student might practice developing it.}

\textit{\textbf{Example}: for the virtue of friendliness, a concrete practice might be asking follow-up questions whenever you meet someone new.}

\vspace{0.2\baselineskip}

\wsgridtable{Virtue}{Concrete practice(s)}{4}{0.85in}

\vspace{0.4\baselineskip}

\wsprompt{According to virtue ethics, a virtue and a vice often exist on the same scale: the virtue is the "golden mean" between having too little of something and too much of it. As a group, choose some example virtues, then work out what they would look like in \textbf{deficiency} and in \textbf{excess}. For example, excess honesty might be "boastful."}

\textit{\textbf{Some virtues to choose from}: courage, ambition, generosity, patience, honesty, whimsy / wittiness, friendliness --- and many more!}

\vspace{0.2\baselineskip}

\wsgridtablethree{Deficiency (vice)}{Virtue (golden mean)}{Excess (vice)}{3}{0.8in}

\end{document}
